% !PRINTED 410
\weberpage{410}
%<<CONT
it follows that $r^\mu = 1$ must hold. Thus no term that has already occurred can recur in
the sequence (6) before the first term 1 occurs for the second time, and if $\mu$ is the
least positive number for which $r^\mu = 1$, then the numbers:

\weq{7}{1, r, r^2 \ldots r^{\mu - 1}}
are all different from one another.

\gesperrt{Then $\mu$ must be a divisor of $n$.} For by division the integers $h, \mu'$ can be
determined so that

\[
n = h\mu + \mu'; \quad 0 \leqq \mu' < \mu.
\]
But then, as follows from

\[
r^n = r^{h\mu} r^{\mu'}
\]
we also have $r^{\mu'} = 1$, that is, since $\mu$ is to be the least positive number for
which $r^\mu = 1$, $\mu' = 0$, or $n$ is divisible by $\mu$.

\gesperrt{Thus $r$ is at the same time a $\mu$th root of unity, but not a root of unity of any
still lower degree.}

\gesperrt{The number $r$ is called a primitive $n$th root of unity if it is not at the
same time a root of unity of a lower degree.}

From this definition it follows that the numbers of the sequence (6), if $r$ is a
primitive $n$th root of unity, are all different from one another, and that all the
$n$th roots of unity are contained among them, but that they make up only a part of the
$n$th roots of unity if $r$ is an imprimitive $n$th root of unity.

Every root of unity whose degree is a divisor of $n$ different from $n$ is at the same
time an imprimitive $n$th root of unity.

\gesperrt{If $r$ is at the same time an $n$th and an $m$th root of unity, then it is also a
$\mu$th root of unity, if $\mu$ is the greatest common divisor of $n$ and $m$.}

For by \S.\,118 the integers $x, y$ can be determined so that

\[
m x + n y = \mu
\]
and consequently

\[
r^\mu = r^{mx} r^{ny} = 1.
\]
To this corresponds the other theorem:

\gesperrt{If $r_1, r_2 \ldots$ are roots of unity of the degrees $n_1, n_2 \ldots$, then they are
all at the same time roots of unity of the degree $m$, if $m$ denotes any common
multiple of $n_1, n_2 \ldots$.}
